\section{13. Verification as an Assurance Argument}

Sections 4 and 7 already establish that verification is not a single binary property and that "Oracle strength" varies by task. This section makes that idea fully operational: for any specific check, "verified" is an incomplete sentence until it answers \textbf{verified against what, for which failure mode, with what coverage and independence, and under what conditions?}

Six properties are worth recording for any verification mechanism actually in use:

\begin{center}
\begin{tabular}{|p{0.22\textwidth}|p{0.68\textwidth}|}
\hline
Property & Question \\
\hline
\textbf{Target} & What specific failure is this check intended to detect? \\
\hline
\textbf{Coverage} & Which claims, steps, or cases does it actually examine? \\
\hline
\textbf{Sensitivity} & How often does it detect the targeted error when the error is present? \\
\hline
\textbf{Specificity} & How often does it incorrectly reject a sound result? \\
\hline
\textbf{Independence} & What assumptions, data, models, or infrastructure does it share with whatever produced the thing it's checking? \\
\hline
\textbf{Freshness} & When does the evidence underlying this check expire? \\
\hline
\end{tabular}
\end{center}

This makes verification gaps visible in a way a simple "was it checked, yes or no" cannot. A source-grounding check, from Section 4, might have excellent sensitivity for catching misquotation while having precisely zero coverage of omitted authority, stale law, or a flawed strategic judgment --- "verified" and "unverified" are both misleading one-word summaries of that situation; the honest description names the specific target the check actually has.

\textbf{Completeness and omission deserve their own explicit treatment, separate from claim checking, and this may be the single most consequential idea in this section.} Most verification methods are considerably better at checking claims that \emph{appear} than at detecting required claims that \emph{never appeared at all} --- the worked example in Section 12 is exactly this failure mode. A citation checker can confirm that every proposition actually made is well-supported while completely missing that the analysis omitted an adverse authority, a contractual exception, a relevant jurisdiction, a document that should have been in scope, an alternative interpretation worth raising, a deadline, or a required approval. Four genuinely separate questions are worth keeping apart: \textbf{claim validity} (are the claims that were made correct?); \textbf{set completeness} (are all the issues that should be represented actually represented?); \textbf{search coverage} (was the relevant universe of material adequately examined in the first place?); and \textbf{scope adequacy} (was that universe itself correctly defined to begin with?). Completeness needs a different kind of verification standard than claim-checking does --- a checklist, a taxonomy, a defined issue universe, a document inventory, a coverage sample, or a qualified reviewer specifically asked "what's missing here" rather than "is this right." And any negative conclusion --- "no relevant clause exists," "no adverse authority was found" --- should carry an explicit coverage statement rather than being stated as an unqualified fact: \emph{"no relevant clause was found in the 37 documents indexed as of this date, searched against these clause categories"} is a defensible, checkable claim; \emph{"there is no relevant clause"} is not, and conflating the two is exactly how an omission becomes invisible.

A related concept worth naming: \textbf{assurance debt.} Every important proposition an output relies on, without an adequate checker actually verifying it, is a specific, nameable, owned obligation --- not something that should be allowed to disappear into a general statement that "human review is required" and then never get tracked again. Treating it as debt (visible, attributable, and expected to eventually be paid down or explicitly accepted) is a meaningfully different posture than treating it as a vague caveat.

\textbf{Consequence mechanics: why a lower error rate is not the same thing as a safer system.} A workflow's raw accuracy doesn't determine how much an error actually matters operationally; a fuller picture needs: \textbf{severity} (how harmful is the resulting outcome); \textbf{detectability} (how likely is the error to be noticed at all); \textbf{detection latency} (how long before it's noticed, if it is); \textbf{reversibility} (can the effect be undone once it's happened); \textbf{blast radius} (how many people, matters, records, or systems are affected); \textbf{persistence} (does the error fade, or does it remain embedded in memory and records); \textbf{recurrence} (can the same root cause repeatedly regenerate the same error); \textbf{observability} (will the organization even know the relevant event occurred); and \textbf{attribution} (can the organization identify which component or decision was actually responsible). This is why a workflow with a \emph{lower} measured error rate can still be \emph{less} safe than one with a higher rate: rare, silent, irreversible errors can dominate total harm even when frequent, visible, easily corrected ones would show up as the larger number in a simple accuracy comparison. A useful risk statement names all of these at once --- \emph{"this failure is uncommon but hard to detect, propagates directly into client-facing advice, and can't be fully reversed once it's been communicated"} --- rather than collapsing them into a single accuracy percentage. This kind of failure-mode-to-consequence linkage, rather than stopping at component-level correctness, is the organizing idea behind FMECA and reliability engineering more broadly, and is worth borrowing rather than reinventing.

\textbf{Base rates and asymmetric error costs close out this section.} A system's average accuracy is close to meaningless without knowing the prevalence of whatever condition it's trying to detect. Any classification or detection workflow needs, alongside its accuracy figure: the base rate of the condition being detected; the false-negative rate specifically; the false-positive rate specifically; the cost associated with each type of error, which are rarely equal; performance on rare but consequential subgroups, not just the aggregate; and an abstention or escalation rate. False positives and false negatives propagate differently and shouldn't be collapsed into one number: a false positive tends to create unnecessary review work, and enough of them can desensitize reviewers and erode trust in the system generally; a false negative tends to pass silently into action, with no corrective step ever triggered. Optimizing one of these can straightforwardly worsen the other. The verification-strength axis from Section 7 becomes considerably more operational once paired with an explicit loss function --- the honest question is not just "how strong is this Oracle" but \textbf{"what error trade-off is acceptable for this specific task, and who actually has the standing to make that call?"} --- a decision that belongs with someone accountable for the outcome, not with whoever happened to set the default threshold in the tooling.
